\scanpage{17}
\[
E_i=[\unsure{-i|i},i]\qquad0\leq i\leq m
\]
\margin{NB $E_m=G$, $E_{\unsure{1|0}}=\varepsilon$}
\[
F_{\unsure{i|j}}=[m-j,m+1+j]\qquad0\leq j\leq m-1\qquad F_m=E_m=G
\]
\inserted{\[
E_i\ast E_j=\sum\alpha_{ij}^h E_h\qquad\alpha_{ij}^h\geq0
\]}
\[
F_i\ast F_j=\sum_0^m\beta_{ij}^h E_h\qquad\beta_{ij}^h\geq0
\]
\[
E_i\ast F_j=\sum\gamma_{ij}^h F_h
\]
\margin{$\varepsilon=\varepsilon^{-1}$\quad\deleted{$\unsure{\varepsilon}$}}
\[
\mathcal A(\uline{g,\deleted{\unsure{\delta}}\varepsilon})=
\left\{\begin{array}{ll}
\dfrac g2\ast\mathcal A(0)&\text{si }\varepsilon=+1\\
\dfrac g2\ast\mathcal B(0)&\text{si }\varepsilon=-1
\end{array}\right.
\]
\[
\mathcal A(\tilde g)\ast\mathcal A(\tilde h)\unsure{\subset|\subseteq}\mathcal A(\tilde g\cdot\tilde h)
\]

% The following portion is written in the opposite orientation.
\deleted{$\unsure{G'}\quad G'\text{ groupe}\quad G'\xrightarrow{\lambda}G$}
\[
\tilde G\text{ gr}\qquad\tilde G\xrightarrow{\rho}G\quad\text{hom.}
\]
\[
i:{}_2G\simeq\operatorname{Ker}\rho
\]
\margin{\deleted{\unsure{f}}\quad
$\begin{aligned}
A&=g\ast\unsure{\check A|\tilde A}=\check A\ast g\\
B&=h\ast\check B=\check B\ast h\\
A\ast B&=g\ast\check A\ast h\ast\check B=\unsure{g\ast\check A\ast\check B\ast h}
\end{aligned}$}
\[
\forall\tilde g\in G\qquad\mathcal A(\tilde g)\subset\mathcal P_f(G),
\]
\margin{On pose $C(\tilde g)$ cône des
$\displaystyle\sum_{A\in\mathcal A(\tilde g)}\alpha_A\unsure{A_\ast|A}$, $\alpha_A\geq0$}

a) $\mathcal A(\tilde g)$ \unsure{tot} ordonné

b) $A\in\mathcal A(\tilde g)\Longrightarrow A=\deleted{\unsure{\rho(\tilde g)}}g\ast\check A$ où $g=\rho(\tilde g)$
\margin{$\forall\tilde g\in\tilde G$}

c) $A\in\mathcal A(\tilde g),\ B\in\mathcal B(\tilde h)\Longrightarrow A\ast B\in C(\tilde g\tilde h)$
\margin{$\forall\tilde g,\tilde h\in\tilde G$}

d) Si $u\in{}_2G$, $\tilde g\in G$, \unsure{$\ldots$}

\deleted{d) $\mathcal A(\tilde g\tilde u^2)=\tilde u\ast\mathcal A(\tilde g)\quad\mathcal A(\unsure{\tilde u\tilde g})=$}
$\mathcal A(i(u)\tilde g)=u\ast\mathcal A(\tilde g)$

d) $\forall\tilde g,\tilde h\in\tilde G\quad\exists s\in G$ tel que $\mathcal A(\tilde g)\cup\unsure{\varepsilon_s|s}\mathcal A(\tilde h)$ tot. ordonné

[i.e. $\bigcup\mathcal A_g$
% archive foliation: 16
